\documentclass[10pt,a4paper]{amsart}
\usepackage[margin=19mm]{geometry}
\usepackage{amsmath,amssymb,mathtools}
\usepackage[scr=rsfs,cal=euler]{mathalfa}
\usepackage{xfrac}
\usepackage[unicode,hidelinks,bookmarks=false]{hyperref}
\newcommand{\ie}{i.e.\ }
\newcommand{\cf}{cf.\ }
\newcommand{\resp}{respectively\ }
\newcommand{\Subsection}{\subsection}
\providecommand{\nobreakdash}{\nobreak\hbox{-}}
\setlength{\emergencystretch}{2em}
\multlinegap0pt
\usepackage{bm}
\usepackage{bbm}
\usepackage{dsfont}
\usepackage{xspace}
\usepackage{cancel}
\usepackage{mathtools}
\usepackage{amsthm}
\makeatletter
\@ifundefined{theorem}{\newtheorem{theorem}{Theorem}}{}
\@ifundefined{proposition}{\newtheorem{proposition}[theorem]{Proposition}}{}
\makeatother
\DeclareMathOperator{\Hess}{Hess}
\DeclareMathOperator{\tr}{tr}
\newcommand{\R}{\mathbb{R}}
\newcommand{\bn}{\mathbf{n}}
\newcommand{\cC}{\mathcal{C}}
\newcommand{\cQ}{\mathcal{Q}}
\newcommand{\eps}{\varepsilon}
\newcommand{\pr}{\partial}
\title[Generic mean curvature flows with cylindrical singularities ]{Generic mean curvature flows with cylindrical singularities I: the normal forms and nondegeneracy}
\author{Ao Sun, Jinxin Xue}
\date{Source: arXiv:2210.00419v3 (CC BY 4.0)}
\begin{document}
\maketitle

\noindent\emph{Source and licence.} Generic mean curvature flows with cylindrical singularities I: the normal forms and nondegeneracy. Version arXiv:2210.00419v3 (CC BY 4.0), distributed under the Creative Commons Attribution 4.0 International licence, as recorded in the arXiv licence metadata for this exact version and shown on the abstract page https://arxiv.org/abs/2210.00419v3. Full rights notice: licenses/arxiv-2210.00419-CC-BY-4.0.txt.\par\smallskip

\noindent\textit{Editorial scope.} Counted unit: the Proposition whose source label is \texttt{PropQ} in the article (source0.tex lines 1785--1809, source label \texttt{PropQ}), delivered together with the article's own complete original proof of it (source0.tex lines 1811--2021, one \texttt{proof} environment of 211 lines). No mathematics is written, repaired or re-derived: statement and proof are the article's own text line for line. Required context retained uncounted: none. Every definition, display and label of those lines is retained unchanged. Omitted from this package and not redistributed: 1--1784, 1810--1810, 2022--2025. Nothing inside a retained excerpt is omitted or altered: every retained body is delivered exactly as the article's lines are written, so no omission receipt of any kind accompanies this package. Citations: the retained excerpts carry no citation command at all, so no bibliography entry of the article is transcribed and no reference is redistributed. Reference adaptation: none was needed and none was made; every reference inside the delivered unit resolves inside it, so no unresolved reference or citation is produced. Printed slips kept literal: no printed word, symbol, hypothesis or number of the counted statement or of its proof was altered. Layout and class: the article's own class is \texttt{amsart} with packages including inputenc, amsfonts, appendix, graphicx, psfrag, epsfig, multirow, caption, subcaption, amssymb, amsmath, amscd, amsthm, verbatim; the wrapper is the lane's pinned \texttt{amsart} editorial wrapper at 10pt on A4, which restates the theorem-like environments the retained text needs and adds the article's own macro definitions that the retained text needs (\texttt{\textbackslash Hess}, \texttt{\textbackslash R}, \texttt{\textbackslash bn}, \texttt{\textbackslash cC}, \texttt{\textbackslash cQ}, \texttt{\textbackslash eps}, \texttt{\textbackslash pr}, \texttt{\textbackslash tr}), with extra wrapper packages (bm, bbm, dsfont, xspace, cancel, mathtools). The delivered numbering is the unit's own. Assets: no retained span contains a figure environment, a graphics-inclusion command, a PDF-page inclusion, a table or a TikZ picture; no image file is copied into this package, the package declares no asset dependency and no image file is redistributed with it. Licence: version arXiv:2210.00419v3 of this article is distributed under the Creative Commons Attribution 4.0 International licence, as recorded in the arXiv licence metadata for this exact version and shown on the abstract page https://arxiv.org/abs/2210.00419v3. A full rights notice is delivered at licenses/arxiv-2210.00419-CC-BY-4.0.txt. Characters: every retained excerpt is pure ASCII, so the delivered engine, XeLaTeX with its default Latin Modern text font, reports no missing character.
	\begin{proposition}\label{PropQ}Let $M_t$ be a rescaled mean curvature flow converging to a cylinder $\cC_{n,k}$ in the $C^\infty_{loc}$ sense as $t\to\infty$. Writing $M_t$ as a normal graph of a function $u$ over $\cC_{n,k}$ within the graphical radius, then we have  
		\begin{enumerate}
			\item 	we have
			$\partial_tu=L u+\cQ(J^2u),\ J^2u:=(u,\nabla u, \nabla^2u),$
			where we have 	\begin{equation}\label{EqQ<}
				|\cQ(J^2u)|\leq 
				C(|\nabla u|^4+|\nabla u|^2|\Hess_u|+|\nabla u|^2+u^2+|u||\Hess_u|)
			\end{equation}
			if $\|u\|_{C^2}\leq \eps_0$ for some $\eps_0$ small. 
			\item 	The leading terms in $\cQ$ is given explicitly as 
			\begin{equation}\label{EqQ}
				\cQ(J^2u)=-(2\varrho)^{-1}\left(u^2+4u\Delta_\theta u+2|\nabla_\theta u|^2\right) +\cC(J^2u),	
			\end{equation}
			where $\cC(J^2u)$ consists of terms cubic and higher power in $J^2u$, satisfying
			\begin{equation}
				|\cC(J^2u)|\leq C(|u|^3+|\nabla u|^3+|\nabla^2u|(u^2+|\nabla u|^2))\leq C\eps_0(u^2+|\nabla u|^2).
			\end{equation} 
			\item For three functions $v, u_1,u_2:\ \Sigma\cap B_R\to \R$ and the cutoff function $\chi$ that is 1 on $B_{R-1}$ and $0$ outside $B_R$, we have $($denoting $w=u_1-u_2)$\begin{equation*}\label{EqdC}
				\int_{\Sigma\cap B_R} |v 	(\cC(J^2(\chi u_1))-	\cC(J^2(\chi u_2)))|e^{-\frac{|x|^2}{4}}\leq C\max\{\|u_1\|_{C^2},\|u_2\|_{C^2}\}^2(\|w\|\cdot \|v\|+e^{-\frac{R^2}{4}}).
			\end{equation*}
		\end{enumerate}
		
		
		
	\end{proposition}

	\begin{proof}
		
		Note that the unit normal vector $\bn=\frac{\theta}{\varrho}$. Later we will use the following fact:
		\[\pr_i \theta=\pr_i\pr_j\theta=0,\ \ \ 
		\pr_\alpha\theta=\theta_\alpha,\ \ \ \pr_{\alpha}\pr_{\beta} \theta = -\frac{\delta_{\alpha\beta}}{\varrho^2}\theta.
		\]
		We consider a graph $\cC_{n,k}$ locally given by $\{\widetilde{F}(z)=(z)+u(z)\bn= (z)+u(z)\frac{\theta}{\varrho}\}$, which induces a frame on $\cC_{n,k}$ given by
		$
		\widetilde{\pr}_\alpha=(1+\frac{u}{\varrho})\theta_{\alpha}+\frac{\pr_\alpha u}{\varrho} \cdot\theta,\ \ \ 
		\widetilde{\pr}_{i}=\pr_i +\frac{\pr_iu}{\varrho} \cdot \theta.
		$
		The induced metric is given by
		\begin{equation*}
			\widetilde{g}_{\alpha\beta}=(1+\frac{u}{\varrho})^2 g_{\alpha\beta} + \pr_\alpha u \pr_\beta u
			,
			\ \ \ 
			\widetilde{g}_{i j}
			=
			g_{ij}+\pr_i u\pr_j u
			,
			\ \ \ 
			\widetilde{g}_{\alpha i}
			=
			\pr_{\alpha} u \pr_i u.
		\end{equation*}
		The inverse matrix is given by
		\begin{equation*}
			\begin{aligned}
				\widetilde{g}^{\alpha\beta}
				&=
				(1+\frac{u}{\varrho})^{-2}g_{\alpha\beta} - (1+\frac{u}{\varrho})^{-4}\pr_\alpha u \pr_\beta u+ m_{\alpha\beta},
				\quad	\\
				\widetilde{g}^{i j}
				&=
				g_{ij} - \pr_i u\pr_j u 
				+ m_{ij},\quad 
				\widetilde{g}^{\alpha i}
				=
				-\pr_\alpha u\pr_i u
				+ m_{\alpha i}.
			\end{aligned}
		\end{equation*}
		Here $m$ are functions of products of $\nabla u$, at least quadratically.	We find a unit normal vector field given by
		\begin{equation*}
			\widetilde{\bn}
			=
			\frac{\frac{\theta}{\varrho}-(1+\frac{u}{\varrho})^{-1}\pr_\alpha u\theta_{\alpha} -\pr_i u \pr_i}
			{\sqrt{
					1+\sum_{\alpha}(1+\frac{u}{\varrho})^{-2}(\pr_\alpha u)^2	
					+\sum_i (\pr_i u)^2.
			}}
			=:
			\frac{\frac{\theta}{\varrho}-(1+\frac{u}{\varrho})^{-1}\pr_\alpha u\theta_{\alpha} -\pr_i u \pr_i}{S}.
		\end{equation*}
		
		Now we calculate the 2nd fundamental form. We have
		\begin{equation*}
			\begin{aligned}
				\pr_{\alpha}\pr_{\beta} F
				&=
				-(1+\frac{u}{\varrho})\frac{\delta_{\alpha\beta}}{\varrho^2}\theta
				+
				\frac{\pr_{\beta} u}{\varrho}\theta_{\alpha}
				+
				\frac{\pr_{\alpha }u}{\varrho}\theta_\beta
				+
				\frac{\pr_{\alpha\beta}u}{\varrho}\theta,
				\\
				\pr_i\pr_{\alpha} F
				&	=
				\frac{\pr_{i\alpha} u}{\varrho}\theta
				+\frac{\pr_i u}{\varrho}\theta_{\alpha},
				\quad		\pr_i\pr_j F
				=
				\frac{\pr_{ij}u}{\varrho}\theta.
			\end{aligned}
		\end{equation*}
		
		Taking inner product with $\widetilde{\bn}$, we get the 2nd fundamental forms:
		\begin{equation*}
			\begin{aligned}
				\tilde{A}_{\alpha\beta}
				&=
				S^{-1}\left(
				-(1+\frac{u}{\varrho})\frac{\delta_{\alpha\beta}}{\varrho}
				-2(1+\frac{u}{\varrho})^{-1}\frac{\pr_\alpha u \pr_{\beta}u}{\varrho}
				+\pr_{\alpha\beta} u
				\right)	,\\
				\tilde{A}_{i\alpha}
				&=
				S^{-1}\left(
				\pr_{i\alpha} u 
				-
				(1+\frac{u}{\varrho})^{-1}\frac{\pr_{\alpha}u \pr_i u}{\varrho}
				\right)	,
				\quad 		\widetilde{A}_{ij}
				=
				S^{-1}\left(\pr_{ij}u
				\right)	.
			\end{aligned}	
		\end{equation*}
		
		In conclusion, we have (note that the convention in mean curvature flow is $H=-\tr A$)
		\begin{equation}\label{eq:H}
			\begin{split}
				-\widetilde{H}
				=&
				S^{-1}
				\left(
				(1+\frac{u}{\varrho})^{-2}\Delta_\theta u
				+\Delta_{\R^k} u
				-\pr_{ij}u\pr_i u\pr_j u
				-\pr_{i\alpha} u\pr_i u\pr_\alpha u
				\right.
				\\
				&
				\left.-(1+\frac{u}{\varrho})^{-4}\pr_\alpha u \pr_\beta u\pr_{\alpha\beta}u-(1+\frac{u}{\varrho})^{-1}\frac{k}{\varrho}
				-(1+\frac{u}{\varrho})^{-3}\frac{|\nabla_\theta u|^2}{\varrho}
				+m
				\right)
			\end{split}
		\end{equation}
		where $m$ is a remainder function of the linear combinations of products of $\nabla u$, with order at least $4$. In fact, we have
		\begin{equation}\label{Eqm1}
			|m|\leq C|\nabla u|^4,\quad  |\nabla m|\leq C|\nabla u|^3.
		\end{equation}
		
		Next we consider the term $\frac{1}{2}\langle \widetilde z,\widetilde \bn\rangle$. On $\cC_{n,k}$ near $z_0=(\theta_0,y_0)$, we have
		\begin{equation*}
			\begin{split}
				\langle \widetilde z,\widetilde\bn\rangle
				=&
				S^{-1}\left\langle (z)+\frac{u}{\varrho}\theta
				, \frac{\theta}{\varrho}-(1+\frac{u}{\varrho})^{-1}\pr_\alpha u\theta_{\alpha} -\pr_i u \pr_i
				\right\rangle
				=
				S^{-1}
				\left(
				\varrho+u
				-y_i\pr_i u
				\right)
			\end{split}
		\end{equation*}
		Therefore, we can obtain the equation of $u$ from the rescaled mean curvature flow equation (modulo diffeomorphism) $(\pr_t \widetilde z)^{\bot}=-(\widetilde H-\frac{\langle \widetilde z,\widetilde\bn\rangle}{2})\widetilde \bn$. Take an inner product of the equation with $\widetilde \bn$, we get 
		\begin{equation*}
			\begin{split}
				S^{-1} \pr_t u
				=&
				S^{-1}
				\left(
				(1+\frac{u}{\varrho})^{-2}\Delta_\theta u
				+\Delta_{\R^k} u
				-\pr_{ij}u\pr_i u\pr_j u
				-\pr_{i\alpha} u\pr_i u\pr_\alpha u
				-(1+\frac{u}{\varrho})^{-4}\pr_\alpha u \pr_\beta u\pr_{\alpha\beta}u
				\right.
				\\
				&
				\left.-(1+\frac{u}{\varrho})^{-1}\frac{k}{\varrho}
				-(1+\frac{u}{\varrho})^{-3}\frac{|\nabla_\theta u|^2}{\varrho}
				+m
				\right)+S^{-1}\frac{1}{2}(\varrho+u-y_i\pr_i u).
			\end{split}
		\end{equation*}
		Thus we obtain the equation of $u$ as follows:
		\begin{equation}\label{EqRMCFFull}
			\begin{split}
				&\pr_t u
				=
				\left(
				(1+\frac{u}{\varrho})^{-2}\Delta_\theta u
				+\Delta_{\R^k} u
				-\pr_{ij}u\pr_i u\pr_j u
				-\pr_{i\alpha} u\pr_i u\pr_\alpha u	-(1+\frac{u}{\varrho})^{-1}\frac{n-k}{\varrho}\right.\\
				&\left.-(1+\frac{u}{\varrho})^{-4}\pr_\alpha u \pr_\beta u\pr_{\alpha\beta}u
				-(1+\frac{u}{\varrho})^{-3}\frac{|\nabla_\theta u|^2}{\varrho}
				+m
				\right)
				+\frac{1}{2}(\varrho+u-y_i\pr_i u).
			\end{split}
		\end{equation}
		When $|u|_{C^0}\leq \eps_0$, we can further rewrite the equation as 
		$\pr_t u
		=
		Lu+\cQ(J^2u),$
		where $\cQ(J^2u)=m_1+m_2+m_3+m_4$, with $m_4$ consisting of all terms like $\pr u\pr u\pr^2 u$ and 
		\begin{equation*}
			m_1=m,
			\ \ \ m_2=-\frac{1}{2\varrho}\frac{u^2}{1+\frac{u}{\varrho}}-(1+\frac{u}{\varrho})^{-3}\frac{|\nabla_\theta u|^2}{\varrho},
			\ \ \ 
			m_3=-\frac{\frac{2u}{\varrho}+\frac{u^2}{\varrho^2}}{(1+\frac{u}{\varrho})^2}\Delta_\theta u.
		\end{equation*}
		From \eqref{EqRMCFFull}, we get \eqref{EqQ} immediately. Indeed, the leading terms come from $m_2$ and $m_3$ respectively. The terms $m_1$ (see \eqref{Eqm1}) and $m_4$ contribute only to $\cC(J^2u)$. Estimate \eqref{EqQ<} follows also immediately from \eqref{EqRMCFFull}. In particular, when $\|u\|_{C^{2}}\leq \eps_0$, we have the estimate
		$		|\cQ|\leq C(|\nabla u|^2+\eps_0u).
		$	
		For two different functions $u_1$ and $u_2$, the fundamental theorem of calculus shows that
		\begin{equation*}
			|\cQ(J^2u_1)-\cQ(J^2u_2)|\leq C\eps_0(|u_1-u_2|+|\nabla u_1-\nabla u_2|+|\Hess_{u_1-u_2}|).
		\end{equation*}
		We next consider item (3). To get the terms $\|u\|\|v\|$ on the RHS, we need to perform a step of integration by parts for terms of the form $v\mathrm{Hess}_u$, which also gives us the boundary term $e^{-R^2/4}$. Taking terms $\partial_iu\partial_j u\partial_{ij}u$ in $m_4$ as an example, we have $$\partial_iu\partial_j u\partial_{ij}u=(\nabla u)^T\nabla^2u\nabla u=\frac{1}{2}\nabla\cdot (|\nabla u|^2\nabla u)-\frac{1}{2}\Delta u \cdot |\nabla u|^2. $$ 
		Consider the first term on the RHS
		\begin{equation*}
			\begin{aligned}
				&\nabla\cdot (|\nabla (\chi u_1)|^2\nabla (\chi u_1))-\nabla\cdot (|\nabla (\chi u_2)|^2\nabla (\chi u_2))\\
				&=\nabla\cdot ((\nabla (\chi w)\cdot\nabla(\chi u_1+\chi u_2)) \nabla (\chi u_1))+\nabla\cdot (|\nabla (\chi u_2)|^2\nabla (\chi w))). 		
			\end{aligned}
		\end{equation*}
		When multiplied by $v$ and taking integration by parts, we see that it can be bounded by the RHS of \eqref{EqdC}. 
		When the derivative during the integration by parts hits the Gaussian weight $e^{-\frac{|x|^2}{4}}$, we shall get $\int v|x|e^{-\frac{|x|^2}{4}}$ which is bounded by $\|v\|_{H^1}$ using the Ecker's inequality $\int v^2|x|^2e^{-\frac{|x|^2}{4}}\leq C\int (v^2+|\nabla v|^2) e^{-\frac{|x|^2}{4}}. $ 
		All other terms can be treated similarly and are easier, so we get \eqref{EqdC}. 
	\end{proof}

\end{document}
